\chapter{Dostępność i~użyteczność}\label{ch:dostepnosc}

\section{Zasady przyjęte w~projekcie}

Dostępność potraktowano jako wymaganie nadrzędne -- dotyczy każdej strony platformy,
również stron logowania przez portale, panelu administracyjnego i~komunikatów
systemowych. Interfejs nie korzysta z~żadnej biblioteki komponentów; każdy
element interaktywny został napisany od podstaw w~HTML, CSS i~JavaScripcie, z~semantyką
rodzimych elementów (\code{button}, \code{a}, \code{details}, \code{dialog},
\code{fieldset}) zamiast \code{div} z~obsługą kliknięć.

\section{Zgodność z~kryteriami WCAG 2.1}

{\small
\begin{longtable}{@{}p{4.3cm}p{10.3cm}@{}}
\caption{Realizacja wybranych kryteriów WCAG 2.1 (poziom A i~AA)}\label{tab:wcag}\\
\toprule
\hd{Kryterium} & \hd{Realizacja} \\
\midrule
\endfirsthead
\toprule
\hd{Kryterium} & \hd{Realizacja} \\
\midrule
\endhead
\bottomrule
\endlastfoot
1.1.1 Treść nietekstowa & wymagany tekst alternatywny dla zdjęć galerii, zdjęć w~pytaniach i~obrazów w~edytorze lekcji; ikony dekoracyjne z~\code{aria-hidden} \\
1.3.1 Informacje i~relacje & landmarki \code{header}/\code{nav}/\code{main}/\code{footer}, jeden \code{h1} na stronie, logiczna hierarchia nagłówków, tabele z~\code{th}, grupy pól w~\code{fieldset}/\code{legend} \\
1.4.1 Użycie koloru & statusy, błędy i~stany (np.\ ,,ukończona'', ,,obowiązkowy'') zawsze z~tekstem lub ikoną, nie tylko kolorem \\
1.4.3 Kontrast (minimum) & kolor tekstu wyliczany z~kontrastu dla każdego tła i~każdego motywu (rozdz.~9.1) \\
1.4.4 Zmiana rozmiaru tekstu, 1.4.10 Dopasowanie do ekranu & jednostki \code{rem}/\%, układ mobile-first; brak poziomego przewijania przy 320 px i~zoomie 200\% \\
1.4.11 Kontrast elementów nietekstowych & ramki pól i~wskaźnik fokusu o~kontraście $\geq$ 3:1 \\
1.4.13 Treść spod kursora lub fokusu & panele (powiadomienia, menu konta, motyw) otwierane świadomie, zamykane Esc \\
2.1.1 Klawiatura, 2.1.2 Bez pułapki & wszystkie funkcje z~klawiatury, w~tym DnD (przyciski $\uparrow\downarrow$) i~slajder ($\leftarrow\rightarrow$); fokus nie utyka w~żadnym elemencie \\
2.2.2 Pauza, zatrzymanie & slajder z~przyciskiem pauzy, zatrzymywany przy fokusie i~najechaniu; bez automatu przy \code{prefers-reduced-motion} \\
2.4.1 Pomijanie bloków & link ,,Przejdź do treści'' jako pierwszy element strony \\
2.4.3 Kolejność fokusu & zgodna z~kolejnością wizualną, brak dodatnich \code{tabindex} \\
2.4.7 Widoczny fokus & \code{:focus-visible} z~obrysem w~kolorze akcentu na każdym elemencie, nigdy \code{outline: none} bez zamiennika \\
2.4.8 Lokalizacja & okruszki nawigacyjne, \code{aria-current="page"} w~menu i~paginacji \\
2.5.5 Rozmiar celu & przyciski i~linki menu co najmniej 44$\times$44 px \\
3.1.1 Język strony & \code{<html lang="pl">} \\
3.2.2 Przy wprowadzaniu danych & filtry katalogu nie wysyłają formularza same -- wymagają przycisku \\
3.3.1 Identyfikacja błędów, 3.3.3 Sugestie & komunikat tekstowy przy polu, \code{aria-invalid}, \code{aria-describedby}; zbiorczy komunikat \code{role="alert"} nad formularzem \\
3.3.2 Etykiety lub instrukcje & każde pole ma \code{<label for>}; placeholder nigdy nie zastępuje etykiety \\
4.1.2 Nazwa, rola, wartość & \code{aria-expanded}, \code{aria-pressed} (głosy), \code{aria-controls}; własne kontrolki zbudowane na natywnych elementach \\
4.1.3 Komunikaty o~stanie & toasty, wynik zapisu kolejności, licznik koszyka i~gotowa odpowiedź asystenta w~regionach \code{aria-live}/\code{role="status"} \\
\end{longtable}
}

\section{Obsługa klawiaturą}

\begin{table}[H]
\centering\small
\caption{Skróty i~zachowania klawiatury w~komponentach}
\begin{tabularx}{\linewidth}{@{}P{3.8cm}P{3.6cm}L@{}}
\toprule
\hd{Komponent} & \hd{Klawisze} & \hd{Zachowanie} \\
\midrule
Cała strona & Tab, Shift+Tab & przejście między elementami; pierwszy -- ,,Przejdź do treści'' \\
Przełącznik motywu, menu konta & Enter/Spacja, $\uparrow\downarrow$, Home, End, Esc & otwarcie, ruch po pozycjach z~zawijaniem, zamknięcie i~powrót fokusu na przycisk \\
Panel powiadomień & Enter, Tab, Esc & Esc oraz wyjście Tabem poza panel zamyka go \\
Dialog potwierdzenia & Tab (cyklicznie), Esc & jedyny dozwolony ,,trap'' -- Tab krąży w~otwartym modalu; Esc zamyka, fokus wraca do wywołującego \\
Okno asystenta & Enter, Shift+Enter, Esc & wysłanie, nowy wiersz, zamknięcie; panel niemodalny -- Tab może z~niego wyjść \\
Slajder & $\leftarrow\rightarrow$, Spacja/Enter na pauzie & przesunięcie o~jedną kartę, pauza/wznowienie \\
Konfigurator slajdera & Tab, Enter na $\uparrow\downarrow$ / Dodaj / Usuń & alternatywa dla przeciągania; po przeniesieniu fokus zostaje na karcie \\
Program kursu w~kreatorze & Enter na $\uparrow\downarrow$, Esc & zmiana kolejności; edycja nazwy w~\code{<details>} zamykanym Esc \\
Formularze & Enter, Spacja & wysłanie formularza, zaznaczenie pola wyboru \\
Wybór stylu uczenia się & Tab, $\uparrow\downarrow$, Enter & karty stylów to grupa pól radio w~\code{fieldset}; strzałki zmieniają styl, Enter zapisuje \\
Ankieta stylu & Tab, $\uparrow\downarrow$, Enter & każde pytanie to \code{fieldset} z~\code{legend}; strzałki zmieniają odpowiedź w~pytaniu, Tab przechodzi do następnego \\
Wersje lekcji & Tab, Enter & przełącznik wersji to nawigacja z~linkami, bieżąca wersja ma \code{aria-current} \\
Komentarze & Enter na ,,Odpowiedz'' / ,,Edytuj'', Esc & rozwinięcie formularza w~miejscu, zamknięcie Esc \\
\bottomrule
\end{tabularx}
\end{table}

\section{Wspólne komponenty dostępne}

\begin{itemize}
  \item \textbf{Dialog potwierdzenia} (\code{showConfirmDialog}, atrybut
        \code{data-confirm}) zamiast \code{confirm()} -- używany przy usuwaniu, odrzucaniu
        i~anulowaniu. Element \code{<dialog>} z~etykietą, opisem, cyklicznym fokusem i~Esc.
  \item \textbf{Toasty} dla komunikatów \code{django.contrib.messages} -- w~regionie
        \code{aria-live}, nie zabierają fokusu, mają przycisk zamknięcia.
  \item \textbf{Paginacja} (\code{includes/pagination.html}) -- \code{nav} z~etykietą,
        \code{aria-current} na bieżącej stronie, odmiana ,,strona 1 z~5''.
  \item \textbf{Pusty stan} (\code{includes/empty\_state.html}) -- tekst wyjaśniający
        i~przycisk następnej akcji zamiast pustej tabeli.
  \item \textbf{Podgląd zdjęcia przed wysłaniem} -- nazwa, rozmiar, wymiary, ostrzeżenie
        o~limicie, ,,Usuń wybór'' z~powrotem fokusu do pola, komunikaty \code{aria-live}.
\end{itemize}

\section{Responsywność}

Układ projektowano od najmniejszego ekranu. Progi: \textbf{48 rem} (tablet -- siatki
kart w~dwóch kolumnach, slajder 2 karty), \textbf{64 rem} (desktop -- menu poziome,
3 kolumny, okno asystenta jako panel boczny), \textbf{96 rem} (szerokie ekrany -- etykiety
przy ikonach w~nagłówku). Szerokości nagłówka i~okna asystenta zmierzono w~Chrome
dla 320, 768, 1024, 1280, 1366 i~1920 px -- bez poziomego przewijania i~ucinania tekstu.

\section{Unikalny charakter interfejsu}

Język wizualny nawiązuje do \emph{planu zajęć i~zeszytu}: linie zamiast cieni, kanciaste
przyciski, duża typografia (Bricolage Grotesque w~nagłówkach, Atkinson Hyperlegible -- krój
zaprojektowany dla osób słabowidzących -- w~treści), numerowane sekcje 01--04 i~prawdziwe
dane z~bazy zamiast ozdobnych statystyk. Świadomie zrezygnowano z~gradientów, zaokrąglonych
,,pigułek'' i~kart z~ikonami w~kółkach. Ten sam styl obejmuje strony allauth, dialogi,
toasty i~panel administratora (django-unfold w~kolorach motywu).
